% !Mode:: "TeX:UTF-8"
\chapter{绪论}
\label{cha:intro}

\section{研究背景与意义}

\subsection{三维点云数据应用场景}
三维点云是三维空间中离散采样点的集合，能够直接、精准地还原现实世界中物体与场景的三维几何结构，是当前三维视觉领域最核心的数据表示形式之一。与二维图像相比，点云数据不受光照、视角变化带来的尺度歧义影响，能够直接提供物体的三维空间坐标信息，因此在众多工业与科研领域得到了广泛的应用。

在自动驾驶领域，激光雷达（LiDAR）采集的点云数据是环境感知系统的核心数据源，其能够在复杂光照与天气条件下，提供高精度的三维空间位置信息，是自动驾驶车辆实现障碍物检测、可通行区域分割、实时定位与地图构建（SLAM）的关键基础\cite{zhai2023improved}。通过对车辆周围环境的点云数据进行处理，自动驾驶系统能够准确识别被遮挡的车辆、行人等交通参与者，为路径规划与避障决策提供可靠的环境信息，直接决定了自动驾驶系统的安全性与可靠性。

在工业逆向工程与质量检测领域，通过三维扫描设备获取实体零件的点云数据，能够快速重建出零件的高精度三维数字模型，实现产品的数字化设计、逆向仿制与缺陷检测。对于航空航天、汽车制造等高端制造领域，复杂曲面零件的加工精度要求极高，通过对加工后的零件进行点云扫描，与设计的标准模型进行比对，能够快速、精准地检测出加工误差，大幅提升产品质量检测的效率与精度，缩短产品的研发与生产周期。

在增强现实（AR）与虚拟现实（VR）领域，点云数据能够真实还原现实场景的三维结构，为虚拟物体与现实场景的精准融合提供空间基准，提升沉浸式交互体验。通过对现实场景进行三维点云扫描，能够快速构建出场景的高精度三维数字孪生模型，实现虚拟内容与现实场景的几何匹配，让虚拟物体能够与现实环境产生真实的物理交互，广泛应用于数字文旅、虚拟直播、工业远程运维等场景。

在智慧农业领域，通过深度相机获取的作物与果实点云数据，可实现果实的精准定位、长势监测与产量预估，辅助采摘机器人克服枝叶遮挡带来的感知难题，提升自动化作业的精度与效率\cite{zhai2023improved}。针对农业场景中果实被枝叶大面积遮挡的问题，通过点云补全技术恢复果实的完整几何结构，能够让采摘机器人准确获取果实的位置与形状信息，大幅提升采摘的成功率，推动农业生产的自动化与智能化发展。

\subsection{点云数据残缺问题与挑战}
在实际的工程应用中，三维点云数据的获取主要依赖激光雷达、RGB-D相机、结构光三维扫描仪等设备，受限于设备硬件性能、采集环境与物体自身特性，原始采集得到的点云数据往往无法完整还原物体的全部几何结构，普遍存在稀疏、非均匀、残缺的问题\cite{yu2021pointr}，这也是制约后续三维处理任务精度的核心瓶颈。

采集设备的硬件性能是导致点云质量下降的基础因素。激光雷达的角分辨率直接决定了点云的采样密度，对于远距离的物体，采集到的点云往往较为稀疏，无法还原物体的精细几何结构；同时激光雷达的测量范围有限，对于超出量程的区域，无法获取有效的深度信息。RGB-D相机通过红外结构光或飞行时间原理测量深度，对于高反光、透明材质的物体表面，深度测量会出现严重的噪声与缺失；同时受限于相机的视场角，单张深度图像只能覆盖有限的场景范围，无法获取完整的环境信息。结构光扫描仪虽然能够获取高精度的稠密点云，但其对采集环境的光照条件要求较高，同时对于表面纹理单一的物体，容易出现匹配失效，导致点云出现大面积缺失。

物体自身的遮挡效应是导致点云残缺的最核心因素。在点云采集过程中，物体的自身结构、场景中的其他物体都会对采集视点形成遮挡，被遮挡的区域无法被扫描设备捕捉，从而形成大面积的点云缺失\cite{yu2021pointr}。例如在自动驾驶场景中，前方的大型车辆会遮挡后方的小型车辆与行人，导致激光雷达无法获取被遮挡目标的点云数据；在工业零件扫描过程中，零件的复杂内腔、凹槽结构会形成自遮挡，无法被外部的扫描设备捕捉到完整的几何结构。这种遮挡导致的点云缺失，往往会破坏物体的拓扑结构，导致后续的特征提取与处理出现严重的误差。

采集视点的单一性同样会导致点云数据的不完整。单视点采集只能获取物体朝向相机一侧的表面信息，无法获取物体背面、侧面与内部的几何结构，形成大规模的缺失区域。虽然可以通过多视点采集与点云配准，获取物体更完整的点云数据，但在很多实际场景中，受限于采集环境与操作空间，无法实现多视点的环绕采集，只能获取单视点的残缺点云数据。例如在自动驾驶场景中，车辆只能获取自身行驶方向上的环境点云，无法对周围环境进行全视角的环绕扫描；在逆向工程中，对于大型的工业设备，无法实现无死角的多视点采集，只能获取有限视角下的点云数据。

上述因素导致实际获取的原始点云数据，普遍呈现出稀疏、非均匀、不完整的特性，这种数据缺陷不仅破坏了物体原本的几何拓扑结构，还会直接导致后续的目标检测、语义分割、点云配准、三维重建等任务出现严重的精度下降，甚至无法完成正常的处理流程\cite{yan2022shapeformer}。例如在三维重建任务中，残缺的点云会导致重建的表面出现大量的孔洞，无法生成完整的三维网格模型；在目标检测任务中，残缺的点云会导致网络无法准确提取目标的几何特征，出现漏检、误检等问题。因此，如何从残缺的点云数据中，恢复出物体完整的三维几何结构，是三维视觉领域亟待解决的基础性问题。

\subsection{点云补全与三维重建研究价值}
点云补全（Point Cloud Completion）作为三维视觉领域的一项基础性任务，其核心目标是从残缺的输入点云中，推断并恢复出物体完整的三维几何形状，生成稠密、均匀、完整的点云数据。高质量的点云补全算法，能够有效解决实际采集过程中点云残缺的问题，为后续的三维处理任务提供高质量的完整点云数据，因此具备重要的理论意义与实际应用价值。

在理论研究层面，点云补全任务面临着诸多核心的理论挑战，其研究能够推动三维视觉领域相关理论的发展。点云数据的无序性、稀疏性与不规则性，给特征提取与几何建模带来了极大的难度；而当输入点云存在大面积缺失时，如何仅从有限的观测数据中，准确推断出物体缺失区域的几何结构，本质上是一个高难度的逆问题，需要网络具备强大的几何先验学习能力与全局结构推理能力。针对这些问题开展研究，能够进一步深化对三维几何数据深度学习建模的理论认知，推动无序点云特征提取、三维形状生成、长距离几何依赖建模等相关技术的发展，为三维视觉领域的其他任务提供理论与技术支撑。

在实际应用层面，点云补全技术能够直接解决众多工业场景中的实际痛点，具备广泛的应用前景。在自动驾驶领域，通过点云补全技术恢复被遮挡的交通参与者的完整点云，能够大幅提升环境感知系统的检测精度与鲁棒性，让自动驾驶系统能够提前识别被遮挡的风险目标，提升自动驾驶的安全性。在工业逆向工程领域，通过补全扫描得到的残缺点云，能够快速重建出零件完整的三维数字模型，无需进行复杂的多视点扫描与配准，大幅提升逆向设计的效率。在智慧农业领域，通过补全被枝叶遮挡的果实点云，能够让采摘机器人准确获取果实的完整形状与位置信息，提升自动化采摘的成功率。在文物数字化保护领域，针对破损的文物点云进行补全，能够还原文物的原始完整形态，为文物的修复、数字化存档与展示提供支撑。

点云补全是三维重建全流程中的核心前置环节，补全后点云的质量直接决定了三维重建的效果。三维重建的核心目标是从离散的点云数据中，重建出物体连续的表面网格模型，实现物体的数字化表示。当输入点云存在残缺时，传统的三维重建算法无法生成完整、封闭的网格模型，会出现大量的孔洞与结构失真。通过点云补全技术，先将残缺的点云恢复为完整的稠密点云，再进行三维表面重建，能够得到完整、光滑、高精度的三维网格模型，实现从残缺点云输入到完整三维模型输出的全流程处理，推动三维重建技术在更多实际场景中的落地应用。

\section{国内外研究现状}

\subsection{点云补全的传统几何方法与深度学习方法}
点云补全技术的发展，经历了从传统几何方法到深度学习方法的演变过程。早期的点云补全方法主要基于传统的几何理论与数值计算方法，依赖于物体的几何先验知识实现残缺区域的补全。这类方法主要可以分为基于插值的方法、基于模板匹配的方法与基于隐函数的方法三大类。

基于插值的补全方法，核心思想是对残缺区域周围的完整点云进行曲面拟合，利用拟合得到的连续曲面对缺失区域进行插值补全。常用的曲面拟合方法包括径向基函数插值、移动最小二乘曲面拟合、B样条曲面拟合等。这类方法对于小范围的局部缺失，能够得到较为平滑的补全结果，计算过程简单，无需额外的先验数据。但当缺失区域较大、物体几何结构较为复杂时，拟合得到的曲面无法准确还原物体的真实拓扑结构，容易出现结构失真、细节丢失的问题，仅适用于拓扑结构简单、缺失范围较小的场景。

基于模板匹配的补全方法，核心思想是预先构建包含大量完整三维模型的模板库，将输入的残缺点云与模板库中的完整模型进行刚性配准，找到与输入点云匹配度最高的模板模型，再利用模板模型的完整几何结构，对输入点云的缺失区域进行补全。这类方法能够处理较大范围的点云缺失，能够保证补全结果的拓扑结构合理性。但这类方法的补全效果高度依赖模板库的覆盖范围，对于模板库中不存在的物体类别，补全效果会急剧下降，泛化能力极差，无法应用于开放场景中的未知物体补全任务。

基于隐函数的补全方法，核心思想是将物体的表面定义为三维空间中隐式标量函数的等值面，通过对输入的残缺点云进行求解，得到能够表征物体完整形状的隐函数，再通过隐函数实现点云的补全与表面重建。典型的方法包括泊松重建、水平集方法等，这类方法能够生成封闭、光滑的物体表面，对噪声具有一定的鲁棒性。但这类方法对输入点云的完整性与法向量精度要求较高，当输入点云存在大面积缺失时，无法准确求解得到符合物体真实结构的隐函数，容易出现严重的结构失真，仅适用于缺失范围较小的点云补全场景\cite{yan2022shapeformer}。

整体而言，传统的点云补全方法依赖于人工设计的几何先验，对于复杂拓扑结构、大面积缺失的点云，补全效果十分有限，无法满足实际应用场景的需求。随着深度学习技术在计算机视觉领域的快速发展，基于数据驱动的深度学习方法，逐渐成为点云补全任务的主流技术路线。这类方法通过在大规模的完整-残缺点云配对数据集上进行训练，自动学习从残缺点云到完整点云的映射关系，无需人工设计几何先验，具备更强的复杂结构建模能力与泛化能力。

PointNet与PointNet++的提出，为无序点云的深度学习处理奠定了核心基础。这两类网络通过共享的多层感知机（MLP）提取每个点的独立特征，再通过最大池化操作聚合得到点云的全局特征，解决了点云无序性带来的处理难题，首次实现了端到端的点云特征提取与分类、分割任务\cite{qi2017pointnet,qi2017pointnetplusplus}。基于PointNet++的核心架构，PCN（Point Completion Network）首次提出了端到端的点云补全编码器-解码器框架，该框架通过PointNet++编码器提取残缺点云的全局特征，再通过解码器结合折叠操作，从全局特征中生成完整的三维点云，实现了端到端的点云补全\cite{fan2017pcn}。

但PCN仍存在明显的不足，其解码器仅采用简单的二维平面折叠操作生成点云，无法有效建模复杂的拓扑结构，难以还原物体的局部几何细节；同时编码器采用最大池化操作聚合全局特征，会丢失大量的局部结构信息，导致补全结果往往只能还原物体的整体轮廓，无法生成高保真的精细表面\cite{yu2021pointr}。后续的大量研究围绕PointNet++架构展开改进，通过引入多尺度特征融合、残差连接、多级生成策略等方式，提升网络的特征提取能力与补全效果。但这类基于CNN与MLP的方法，其核心的特征提取模块感受野范围有限，无法有效建模点云中点与点之间的长距离依赖关系，对于存在大面积缺失的点云，难以准确推断出缺失区域与完整区域之间的几何关联，补全效果仍存在较大的提升空间\cite{chen2024geoformer}。

\subsection{Transformer在点云处理中的应用}
Transformer架构最初被提出用于自然语言处理任务，其核心的自注意力机制能够有效建模序列中元素之间的长距离依赖关系，相较于CNN与MLP，具备更强的全局特征建模能力。随着Vision Transformer（ViT）在二维图像领域取得的巨大成功，学者们开始将Transformer架构引入三维点云处理领域，其强大的全局上下文建模能力，为点云补全任务带来了新的技术突破\cite{yu2023adapointr}。

PoinTr是将Transformer引入点云补全领域的开创性工作之一，该方法针对点云的无序性与三维几何特性，设计了适配点云补全任务的Transformer架构。针对直接将高分辨率点云作为Transformer输入带来的计算量爆炸问题，PoinTr将输入的残缺点云转换为一组点代理（Point Proxies）序列，每个点代理代表了点云的一个局部几何区域，在大幅降低序列长度的同时，充分保留了点云的核心几何信息。在此基础上，该方法设计了几何感知的Transformer Block，通过将点的三维几何坐标融入注意力权重的计算中，让网络能够同时捕捉点云的局部几何结构与全局形状上下文，有效建模了点云局部区域与全局形状之间的关联，显著提升了点云补全的完整性与结构一致性\cite{yu2021pointr}。

在此基础上，AdaPoinTr进一步针对PoinTr中固定点代理数量带来的计算冗余问题，引入了自适应采样策略，能够根据输入点云的几何复杂度，动态调整点代理的数量与分布。对于几何结构简单的区域，采用较少的点代理降低计算量；对于几何结构复杂的细节区域，采用更多的点代理保留精细的几何信息。该方法在降低计算成本的同时，实现了更高效的几何细节恢复，同时对Transformer Block进行了进一步优化，提升了网络对不同形状物体的泛化能力，相关成果也发表在了模式识别领域的顶级期刊上\cite{yu2023adapointr}。

针对点云无序性带来的计算复杂度问题，ShapeFormer提出了一种基于稀疏体素特征的Transformer架构。该方法首先将输入点云转换为稀疏体素表示，再通过矢量量化编码，将三维体素分布转换为离散的一维序列，将点云补全任务转化为自回归序列预测问题。这种设计有效降低了Transformer的计算复杂度，同时通过稀疏体素表示，更好地保留了点云的三维空间结构信息，提升了网络对复杂形状的建模能力与泛化能力\cite{yan2022shapeformer}。

针对现有Transformer方法在特征提取过程中，局部几何结构捕捉不足的问题，GeoFormer设计了专门的几何Transformer模块。该模块通过构建几何感知的邻域注意力机制，同时捕捉点云的局部精细几何结构与全局形状上下文，通过引入点云的法向量、曲率等局部几何特征，调整注意力权重的分布，让网络更加关注对补全任务更重要的几何边缘与细节区域，有效提升了补全结果的细节保真度\cite{chen2024geoformer}。

另一类研究则聚焦于点云生成范式的优化，通过改进解码器的结构，提升高分辨率点云的生成质量。SnowflakeNet提出了一种雪花点反卷积（Snowflake Point Deconvolution, SPD）技术，配合Skip-Transformer模块，模拟雪花分形生长的过程，从粗糙的骨架点逐步分裂生成精细的高密度点云。该方法在每一个分形生成阶段，都通过Skip-Transformer模块，将编码器提取的多尺度局部特征引入解码器中，让网络能够充分利用输入点云的局部几何信息，极大地丰富了补全结果的表面细节，有效解决了传统方法中细节丢失的问题\cite{xiang2021snowflakenet}。

AnchorFormer则从形状表示的角度出发，通过学习基于区域的锚点表示，利用模式匹配的思想预测缺失区域的几何结构。该方法首先从输入的残缺点云中，提取具有判别性的形状锚点，这些锚点能够表征物体的核心拓扑结构；再通过Transformer建模锚点之间的全局关联，基于锚点表示预测完整的点云形状。该方法在保持补全结果与输入点云的结构一致性方面，取得了优异的效果，能够有效避免补全结果出现全局结构失真的问题\cite{chen2023anchorformer}。

国内学者在点云补全领域同样开展了大量的研究工作，取得了丰硕的研究成果。PCTMA-Net提出了一种结合Transformer与多尺度注意力的网络架构，通过在不同分辨率下提取点云的特征，并进行跨尺度的特征融合，有效解决了单一尺度特征提取过程中信息丢失的问题，提升了网络对不同尺度几何结构的建模能力\cite{lin2023pctmanet}。DMF-Net则构建了一个深层多尺度融合网络，通过在编码阶段设计多分支的特征提取模块，充分保留了输入点云不同尺度的局部特征，同时引入了注意力特征融合模块，实现了局部细节特征与全局形状特征的有效融合，提升了补全结果的整体精度\cite{yang2022dmfnet}。

针对单模态点云数据信息不足，在大面积缺失场景下存在形状歧义性的问题，跨模态点云补全成为了该领域新的研究热点。张文彬等提出了融合图像信息的跨模态Transformer算法，利用高分辨率二维图像提供的纹理与轮廓先验信息，辅助点云几何结构的恢复。该方法设计了跨模态注意力融合模块，实现了图像特征与点云特征的深度交互与融合，利用图像的丰富细节信息，弥补点云数据的局部信息缺失，显著提升了复杂场景下的点云补全质量\cite{zhang2023crossmodal}。Zhang等同样针对多模态点云补全任务开展研究，提出了基于Transformer的多模态特征融合框架，实现了图像与点云信息的互补利用，提升了网络在复杂遮挡场景下对缺失区域的推断能力\cite{zhang2022multimodal}。

为了进一步提升补全结果的精度与细节质量，CRA-PCN设计了级联细化注意力机制，通过多阶段的网络结构，逐步优化生成点云的分布与几何结构。每一阶段的网络都基于上一阶段的补全结果，进行进一步的细节优化与点数上采样，通过级联的多阶段细化，实现了从粗到精的渐进式补全，有效提升了补全结果的细节保真度\cite{yi2022crapcn}。针对实际场景中缺乏完整-残缺配对训练数据的问题，基于自监督学习的点云补全方法逐渐受到了学者们的关注。Lu等提出了一种基于单残缺点云多视图增强的自监督点云补全方法，设计了全新的自监督训练信号，无需配对的完整点云数据，仅利用单张残缺点云即可完成网络的训练，同时首次将Mamba架构引入自监督点云补全任务，在合成数据集与真实世界数据集上都取得了当前最优的自监督补全效果\cite{lu2025selfsupervised}。此外，利用物体的对称性先验信息辅助补全也是该领域的重要研究方向，Yan等提出了SymmCompletion方法，通过对称性引导网络学习物体的对称结构，在提升补全结果保真度的同时，保证了补全结果与原始输入的结构一致性，在高保真点云补全任务上取得了优异的效果\cite{yan2025symmcompletion}。

\subsection{研究现状总结与不足}
综合上述国内外研究现状可以看出，基于Transformer的点云补全技术已经取得了显著的进展，学者们围绕Transformer架构优化、生成范式改进、多模态融合、自监督学习等方向开展了大量的研究工作，有效提升了点云补全的精度与效果，推动了该领域的快速发展。但现有方法仍存在一些尚未解决的问题与不足，无法完全满足实际应用场景的需求，也是本文研究的核心切入点。

首先，现有方法在大面积缺失场景下，难以保证补全后点云全局密度的均衡分布，核心问题集中于原始点云残缺区域的补全密度失衡。当输入点云存在局部残缺时，补全模型对残缺区域的点云生成缺乏精准的密度控制。现有方法往往会出现补全结果中，原始点云残缺区域在补全后的点云密度明显低于全局平均点云密度。根本原因在于模型训练过程中，没有对生成点云的密度不均现象进行针对性惩罚。这种残缺区域补全密度偏低的不均匀现象，会进一步影响后续点云处理任务的性能。例如低密度补全区域难以提供足够的几何特征支撑，导致点云分割、表面重建时出现边界模糊、特征提取不充分的问题，甚至会干扰模型对物体整体结构的精准判断。如何在保证补全形状拓扑正确性的前提下，实现补全后点云全局密度的均衡匹配，解决残缺区域补全密度偏低的问题，避免密度失衡与冗余生成，是该领域需要突破的关键技术瓶颈。

此外，现有方法的实际落地能力仍存在不足。现有方法大多只关注点云补全任务本身的指标提升，很少关注补全算法对下流任务的实际影响，限制了算法的实际落地应用。针对这些问题，本文开展相关研究，设计更高精度的基于Transformer的点云补全算法，并搭建完整的三维重建算法作为下流应用场景，验证算法的实际应用价值。

\section{主要研究内容与技术路线}

\subsection{研究内容}
针对现有基于Transformer的点云补全方法存在的不足，本文以提升点云补全精度、细节保真度与计算效率为核心目标，开展基于Transformer的点云补全与三维重建算法研究，主要研究内容围绕四个核心方面展开。

本文首先开展基于Transformer的点云特征提取网络设计，核心目标是构建能够同时捕捉点云局部几何细节与全局形状上下文的编码器架构。针对点云的无序性与三维几何特性，研究将无序点云转换为适合Transformer处理的序列形式的有效方法，通过点代理生成与自适应采样，在降低序列长度、控制计算复杂度的同时，充分保留点云的核心几何信息。同时，为了保留点云的空间结构信息，引入适配三维点云的位置编码方法，将点的三维空间坐标信息注入序列特征中，解决自注意力机制的排列不变性带来的位置信息丢失问题。在此基础上，设计几何感知的自注意力模块，通过引入点云的局部几何特征调整注意力权重的分布，让网络更加关注对补全任务更重要的几何边缘与细节区域，提升编码器对局部几何细节与全局形状特征的提取能力。

其次，本文开展由粗到精的点云生成策略研究，设计适配Transformer架构的解码器结构，解决高分辨率稠密点云生成难度大、细节丢失的问题。直接从全局特征生成高密度完整点云，会导致网络的学习难度大幅上升，同时难以平衡全局结构与局部细节，因此本文采用由粗到精的生成范式，将点云补全过程分为粗粒度形状生成与细粒度细节优化两个阶段。在粗粒度生成阶段，网络从编码器提取的全局特征中，生成低分辨率的完整点云骨架，还原物体的整体几何结构与拓扑关系，保证补全结果的全局结构一致性，避免出现拓扑结构失真的问题。在细粒度优化阶段，基于粗生成的点云骨架，结合编码器提取的多尺度局部特征，通过点云上采样与细节优化模块，逐步生成高分辨率的稠密点云，还原物体表面的精细几何细节，包括尖锐边缘、微小曲面等，避免出现过度平滑的问题。

此外，本文开展点云补全任务的损失函数优化研究，设计混合损失函数，从多个维度约束生成点云的分布均匀性、几何精度与视觉真实感。传统的点云补全方法大多仅采用倒角距离作为损失函数，虽然计算效率高，但无法有效约束生成点云的分布均匀性，容易出现点云聚集、局部密度不均的问题。因此本文在倒角距离损失的基础上，引入推土机距离损失，提升生成点云与真实点云之间的分布匹配精度；引入密度感知损失，约束生成点云的局部密度，保证点云分布的均匀性；引入排斥损失，避免生成的点出现聚集重叠的问题，提升点云的整体质量。通过对不同损失函数的权重进行优化调整，构建混合损失函数，实现对补全过程的多维度约束，平衡全局结构精度与局部细节质量。

最后，本文开展基于补全算法的三维重建全流程实现研究，搭建包含数据预处理、点云补全、三维表面重建、可视化展示的完整系统，实现从残缺点云输入到完整三维网格模型输出的全流程处理。基于Open3D开源库实现点云的预处理与三维表面重建功能，将本文提出的补全算法作为系统的核心模块，实现残缺点云的高精度补全，再通过泊松重建等方法，从补全后的稠密点云中重建出封闭、光滑的三维网格模型。同时开发可视化交互界面，直观展示点云补全与三维重建的效果，验证算法在实际应用场景中的有效性，推动算法的落地应用。

\subsection{技术路线}
本文的整体技术路线围绕“理论分析-算法设计-实验验证-系统实现”的全流程展开，形成完整的研究闭环，确保研究工作的系统性与严谨性。

在研究的初始阶段，系统梳理点云补全与Transformer相关的国内外研究现状，深入分析现有方法的核心优势与存在的不足，明确本文的核心研究目标与需要解决的关键问题。同时，系统梳理点云数据处理、Transformer架构、三维表面重建相关的基础理论，深入理解相关技术的核心原理与适用场景，为后续的算法设计提供坚实的理论支撑。在此基础上，完成整体算法架构的设计，明确编码器、解码器的核心模块组成，以及混合损失函数的设计方案，形成完整的算法设计思路。同时完成深度学习开发环境的配置，以及ShapeNet、Projected-ShapeNet等公开数据集的下载、预处理与数据加载模块的开发，为后续的算法实现与实验验证奠定基础。

接下来的核心阶段是算法的实现与训练优化，首先完成PoinTr等主流基准算法的复现，跑通开源代码，记录其在测试集上的核心评价指标，作为本文算法的对比基线，确保实验对比的公平性与可靠性。在此基础上，基于PyTorch深度学习框架，完成本文提出的网络架构的代码实现，包括几何感知Transformer编码器、由粗到精的解码器、混合损失函数等核心模块，打通从残缺点云输入到完整点云输出的端到端计算链路。完成代码调试后，开始初步的模型训练与调试，验证网络的收敛性，针对训练过程中出现的梯度消失、不收敛等问题，调整网络结构与超参数，确保网络能够正常训练。

在网络能够正常收敛后，进入大规模模型训练与优化阶段，在训练集上进行完整的模型训练，实时监控训练集与验证集的损失曲线，调整学习率、Batch Size、数据增强策略等超参数，避免模型出现过拟合问题。同时，针对补全效果不理想的复杂拓扑结构样本，针对性地调整网络结构，优化注意力模块与生成策略，提升模型对复杂形状的建模能力与泛化能力。通过多轮的训练与优化，得到最终的最优模型，为后续的实验验证提供基础。

在模型训练完成后，进入实验验证与分析阶段，在测试集上对训练好的模型进行全面的性能评估，计算倒角距离、F-Score等核心评价指标，与主流基准算法进行对比实验，验证本文算法的性能优势。同时开展消融实验，逐一验证本文设计的各个核心模块与损失函数的有效性，分析不同模块对算法性能的影响，明确各个模块的贡献。此外，对补全结果进行可视化对比，直观展示本文算法在全局结构一致性与局部细节恢复上的优势，针对不同缺失程度、不同类别的物体，分析算法的补全效果与鲁棒性。

最后，在算法验证完成的基础上，开展三维重建系统的开发与实现。基于Open3D开源库，实现点云去噪、下采样、归一化、法向量估计等预处理功能，以及泊松重建、球旋转算法等三维表面重建功能。将训练好的补全模型集成到系统中，搭建完整的三维重建全流程系统，实现从残缺点云输入到完整三维网格模型输出的全流程处理。同时开发可视化交互界面，支持点云的加载、显示、补全与重建结果的可视化展示，对系统的运行效果进行全面的测试与分析，验证算法的实际应用价值。

\section{论文结构安排}
本文的整体结构分为五个章节，各个章节的内容安排围绕研究的核心逻辑逐步展开，形成完整、严谨的论文体系。

第一章为绪论，主要阐述本文的研究背景与意义，系统梳理点云补全领域的国内外研究现状，总结现有方法存在的核心不足，明确本文的主要研究内容与技术路线，同时对论文的整体结构安排进行说明，让读者能够对本文的研究工作有整体的认知。

第二章为相关理论与基础技术，主要介绍本文研究涉及的核心基础理论与技术，包括点云数据的基础特性与预处理方法、Transformer架构的核心理论、点云补全任务的核心技术框架与常用损失函数，以及三维表面重建的基础方法，为后续的算法设计、实验验证与系统实现提供坚实的理论基础。

第三章为点云补全算法与三维表面重建算法设计，是本文的核心研究章节，详细阐述本文提出的基于Transformer的点云补全算法的整体网络架构，对点云序列化与几何感知自注意力编码、Coarse-to-Fine渐进式生成策略、混合损失函数的设计进行详细的说明与原理分析。同时，本章还介绍了基于泊松方程的三维表面重建算法的原理与实现细节，包括法线估计、泊松重建、密度过滤和Taubin平滑等关键步骤，为补全点云的下游应用提供完整的重建流程。

第四章为实验设计与结果分析，主要介绍本文的实验环境、采用的数据集与评价指标，通过对比实验与消融实验，对本文提出的算法性能进行全面的验证与分析。对比实验在ShapeNet-34数据集上验证本算法与主流基准方法的性能差异，消融实验逐项分析混合损失函数优化、自适应查询生成机制等核心改进的独立贡献；同时在Projected-ShapeNet数据集上模拟现实扫描产生的残缺点云，对补全结果进行可视化展示，直观验证算法的有效性。为了进一步验证本算法对下游应用场景带来的积极影响，本章将补全后的点云应用于三维表面重建任务，通过可视化展示重建成果来直观反映本算法为下游任务带来的提升效果。

第五章为总结与展望，对本文的研究工作与核心成果进行全面的总结，梳理本文解决的核心问题与创新点，同时分析当前研究存在的不足，对该领域未来的研究方向与优化思路进行展望，为后续的研究工作提供参考。
